\documentclass[11pt]{article}
\usepackage[margin=1in]{geometry}
\usepackage{enumitem}
% =============================================================================
% Preambles/header.tex — shared LaTeX/Beamer preamble for this project
%
% Use from a Beamer lecture by prepending:
%     \input{header}
% and compiling with TEXINPUTS=../Preambles:$TEXINPUTS (the /compile-latex
% skill does this automatically).
%
% PALETTE CONTRACT. Color names defined here MUST match the names in
% Quarto/theme-template.scss so Beamer and Quarto renderings use the same
% palette. The scripts/check-palette-sync.sh script enforces this.
%
% To customize: edit the HEX values below AND the matching $variable in
% Quarto/theme-template.scss, then run scripts/check-palette-sync.sh.
% =============================================================================

% -----------------------------------------------------------------------------
% Engine detection + encoding
%
% Our /compile-latex skill uses XeLaTeX, where inputenc is unnecessary and
% can produce encoding warnings. Only load inputenc under pdfTeX.
% -----------------------------------------------------------------------------
\usepackage{iftex}
\ifPDFTeX
  \usepackage[utf8]{inputenc}
\fi

\usepackage{amsmath, amssymb, amsthm}
\usepackage{booktabs}
\usepackage{graphicx}

% -----------------------------------------------------------------------------
% PALETTE — mirrors Quarto/theme-template.scss variable names.
% Load xcolor and define colors BEFORE anything that references them
% (notably hyperref, hypersetup, and the Beamer color assignments).
% -----------------------------------------------------------------------------
\usepackage{xcolor}

\definecolor{primary-blue}{HTML}{012169}      % headings, accents
\definecolor{primary-gold}{HTML}{B9975B}      % emphasis, borders
\definecolor{highlight-yellow}{HTML}{F2A900}  % markers, alerts
\definecolor{light-bg}{HTML}{E8EDF5}          % title / section backgrounds
\definecolor{jet}{HTML}{1A1A1A}               % body text

% Semantic colors (used in diagrams and annotations)
\definecolor{positive}{HTML}{15803D}          % good / observed / identified
\definecolor{negative}{HTML}{B91C1C}          % bad / problematic / bias
\definecolor{neutral}{HTML}{525252}           % reference / context

% Highlight accents (match .hi-* classes in the SCSS)
\definecolor{hi-slate}{HTML}{314F4F}
\definecolor{hi-green}{HTML}{2E8B57}
\definecolor{hi-red}{HTML}{C41E3A}

% Semantic aliases so skill templates and snippets can write
% \textcolor{accent}{...} without hard-coding "primary-blue".
\colorlet{accent}{primary-blue}
\colorlet{accent2}{primary-gold}

% -----------------------------------------------------------------------------
% Hyperref — loaded AFTER xcolor + palette so link colors resolve.
% -----------------------------------------------------------------------------
\usepackage{hyperref}
\hypersetup{colorlinks=true, linkcolor=primary-blue, urlcolor=primary-blue,
            citecolor=primary-blue, pdfborder={0 0 0}}

% -----------------------------------------------------------------------------
% Course code — set once per deck (after \input{header}, before \begin{document})
% via \coursecode{CS401}. Shown in the footer on every slide so a deck's course
% is identifiable without opening the title page. Empty by default.
% -----------------------------------------------------------------------------
\makeatletter
\newcommand{\coursecode}[1]{\gdef\@coursecodetext{#1}}
\gdef\@coursecodetext{}
\makeatother

% -----------------------------------------------------------------------------
% Beamer theme assignments (only apply under Beamer).
%
% We detect Beamer via \@ifclassloaded{beamer}, which is the documented,
% non-internal way. This must be wrapped in \makeatletter/\makeatother
% because \@ifclassloaded contains an @-catcode character.
% -----------------------------------------------------------------------------
\makeatletter
\@ifclassloaded{beamer}{%
  \usefonttheme{professionalfonts}
  \setbeamercolor{structure}{fg=primary-blue}
  \setbeamercolor{frametitle}{fg=primary-blue}
  \setbeamercolor{title}{fg=primary-blue}
  \setbeamercolor{subtitle}{fg=primary-gold}
  \setbeamercolor{author}{fg=primary-blue}
  \setbeamercolor{institute}{fg=neutral}
  \setbeamercolor{date}{fg=primary-gold}
  \setbeamercolor{itemize item}{fg=primary-blue}
  \setbeamercolor{itemize subitem}{fg=primary-gold}
  \setbeamercolor{alerted text}{fg=primary-gold}
  \setbeamercolor{block title}{fg=white, bg=primary-blue}
  \setbeamercolor{block body}{bg=light-bg}
  % Semantic block variants: block=definition/concept (blue), exampleblock=
  % worked example (green/positive), alertblock=key takeaway or caution (gold).
  % Keep to <=2 colored boxes per slide across ALL three variants (INV-7).
  \setbeamercolor{block title example}{fg=white, bg=positive}
  \setbeamercolor{block body example}{bg=light-bg}
  \setbeamercolor{block title alerted}{fg=white, bg=primary-gold}
  \setbeamercolor{block body alerted}{bg=light-bg}

  % Minimal footer: course code (left) + page X / N (right), no navigation clutter
  \setbeamertemplate{navigation symbols}{}
  \setbeamertemplate{footline}{%
    \color{neutral}\footnotesize\hspace{0.7em}\@coursecodetext\hfill
    \insertframenumber/\inserttotalframenumber\hspace{0.7em}\vspace{0.3em}}
}{}
\makeatother

% -----------------------------------------------------------------------------
% TikZ setup — libraries the snippet gallery uses
% -----------------------------------------------------------------------------
\usepackage{tikz}
\usetikzlibrary{arrows.meta, positioning, calc, decorations.pathreplacing,
                fit, shapes.geometric, backgrounds}

% Shared TikZ styles — snippets and hand-written diagrams can pick these up
% via `\begin{tikzpicture}[every node/.style={...}]` or by naming specific
% styles (e.g., `\node[dag-node] (X) at (0,0) {X};`).
\tikzset{
  % Node styles for causal DAGs and similar
  dag-node/.style = {
    draw=primary-blue, fill=light-bg, rounded corners=2pt,
    minimum width=1.2cm, minimum height=0.9cm, align=center, font=\small
  },
  decision-node/.style = {
    draw=primary-gold, fill=white, diamond, aspect=1.8,
    minimum width=1.8cm, minimum height=1.2cm, align=center, font=\small,
    inner sep=1pt
  },
  % Edge styles — observed vs counterfactual
  observed-edge/.style = {-{Latex[length=2.5mm]}, thick, draw=jet},
  counterfactual-edge/.style = {-{Latex[length=2.5mm]}, thick, draw=neutral, dashed},
  confound-edge/.style = {-{Latex[length=2.5mm]}, thick, draw=negative, dashed},
  % Data-point styles for plots
  observed-dot/.style = {fill=primary-blue, draw=primary-blue, circle, inner sep=1.2pt},
  counterfactual-dot/.style = {fill=white, draw=neutral, circle, inner sep=1.2pt}
}

% -----------------------------------------------------------------------------
% Convenience macros
% -----------------------------------------------------------------------------
\newcommand{\muted}[1]{\textcolor{neutral}{#1}}
\newcommand{\key}[1]{\textcolor{primary-gold}{\textbf{#1}}}
\newcommand{\good}[1]{\textcolor{positive}{#1}}
\newcommand{\bad}[1]{\textcolor{negative}{#1}}

% Standout frame macro: full-bleed dark background for section transitions.
% Beamer-only (uses \begin{frame}/\setbeamercolor/\usebeamerfont) -- do not
% call from an article-class file (e.g. Notes/<CODE>/*-notes.tex). \muted,
% \key, \good, \bad, and \coursecode above are plain xcolor/gdef and are
% safe to use from either class; verified when Notes/article-class support
% was added.
% Expand: \transitionslide{Your transition title}
\newcommand{\transitionslide}[1]{%
  \begingroup
  \setbeamercolor{background canvas}{bg=primary-blue}
  \begin{frame}[plain]
    \centering
    \vfill
    {\usebeamerfont{title}\color{white}\Large #1\par}
    \vfill
  \end{frame}
  \endgroup
}

% Section divider frame (Beamer-only), an alternative to \transitionslide with a
% darker two-line style: near-black (jet) background, a small white label line
% above a larger gold title, and a thin gold rule along the bottom edge.
% Expand: \sectiondivider{Section 2}{Pointers}
\newcommand{\sectiondivider}[2]{%
  \begingroup
  \setbeamercolor{background canvas}{bg=jet}
  \begin{frame}[plain]
    \vspace*{3.0cm}
    {\color{white}\ttfamily\large #1\par}
    \vspace{0.5cm}
    {\color{primary-gold}\LARGE #2\par}
    \begin{tikzpicture}[remember picture, overlay]
      \draw[primary-gold, line width=2.5pt]
        ($(current page.south west)+(0,0.1)$) -- ($(current page.south east)+(0,0.1)$);
    \end{tikzpicture}
  \end{frame}
  \endgroup
}

% -----------------------------------------------------------------------------
% Channel watermark — "Concepts That Click" (YouTube).
%
% Article-class files (CompetitiveExam Q/A sets, Notes, Assignments) get a
% light diagonal draft watermark on every page plus a centered footer naming
% the channel. Beamer decks get a subtle TikZ background watermark instead
% (the footline already carries the course code). Centralized here so every
% MCQ PDF is branded without per-file edits.
% -----------------------------------------------------------------------------
\makeatletter
\@ifclassloaded{article}{%
  \usepackage{draftwatermark}
  \SetWatermarkText{Concepts That Click}
  \SetWatermarkScale{1}
  \SetWatermarkFontSize{2cm}
  \SetWatermarkLightness{0.86}
  \SetWatermarkAngle{45}
  \usepackage{fancyhdr}
  \pagestyle{fancy}
  \fancyhf{}
  \fancyfoot[C]{\footnotesize\textcolor{neutral}{Concepts That Click~|~YouTube: \href{https://www.youtube.com/@concepts.thatclick}{@concepts.thatclick}}}
  \renewcommand{\headrulewidth}{0pt}
  \renewcommand{\footrulewidth}{0pt}
  \fancypagestyle{plain}{%
    \fancyhf{}%
    \fancyfoot[C]{\footnotesize\textcolor{neutral}{Concepts That Click~|~YouTube: \href{https://www.youtube.com/@concepts.thatclick}{@concepts.thatclick}}}%
    \renewcommand{\headrulewidth}{0pt}%
    \renewcommand{\footrulewidth}{0pt}%
  }
}{}
\@ifclassloaded{beamer}{%
  \setbeamertemplate{background}{%
    \begin{tikzpicture}[remember picture, overlay]
      \node[rotate=45, opacity=0.055, align=center]
        at (current page.center)
        {\Huge\textcolor{neutral}{Concepts That Click}};
    \end{tikzpicture}%
  }
}{}
\makeatother 
\coursecode{JKSSB}
\renewcommand{\thesection}{11.\arabic{section}}
\newtheorem{example}{Example}[section]
\newenvironment{definitionbox}[1]{%
  \par\noindent\textbf{Definition (#1).}\ \itshape
}{\par\vspace{0.3em}}
\title{Backup Types --- Detailed Lecture Notes}
\author{JKSSB Computer Awareness}
\date{\today}

\begin{document}
\maketitle

\section{Backup, storage and archive}
These three words describe different jobs, and exam questions often test the
distinction. \textbf{Storage} is where your working data lives day to day ---
the primary copy you actually open and edit. A \textbf{backup} is a separate
copy kept so that data can be recovered after loss, damage, corruption or an
attack. An \textbf{archive} is older data retained for the long term, usually
moved out of routine use. A backup exists for recovery; an archive exists for
retention. Neither is simply ``another folder'' on the same disk.

\begin{definitionbox}{Backup}
A separate copy of data kept for the purpose of recovery after loss, damage,
corruption or attack.
\end{definitionbox}

\begin{definitionbox}{Archive}
Older data kept for long-term retention, normally not used in everyday work.
\end{definitionbox}

The office desktop of the running example (KB EX-01) has one primary copy of
its documents. If the only copy is the one on the desktop's own drive, then
the user has storage but no backup. The moment a second, independent copy
exists, recovery becomes possible.

\section{Full, incremental and differential backups}
A backup type decides \emph{how much} is copied on each run. The three basic
types differ only in the answer to one question: since which earlier point are
changes included?

\begin{definitionbox}{Full backup}
Copies every selected file and folder (or the entire volume) on every run.
It is the recovery baseline for the other types.
\end{definitionbox}

\begin{definitionbox}{Incremental backup}
Copies only the files changed or added since the \emph{last backup of any
type} --- whether that was a full, an incremental or a differential backup.
\end{definitionbox}

\begin{definitionbox}{Differential backup}
Copies everything changed or added since the \emph{last full backup}, not
since the previous backup.
\end{definitionbox}

The single phrase that separates incremental from differential is therefore
the word \emph{full}. An incremental looks only as far back as the
immediately preceding backup; a differential always looks back to the last
full backup, so it keeps growing until the next full backup resets it.

\begin{center}
\begin{tabular}{p{0.20\textwidth}p{0.34\textwidth}p{0.36\textwidth}}
\toprule
\textbf{Type} & \textbf{What it copies} & \textbf{Restore requirement}\\
\midrule
Full & Everything selected, every run. & One complete set. Fastest single restore, but largest and slowest to create.\\
Incremental & Changes since the last backup of any type. & The last full backup plus every incremental backup in sequence.\\
Differential & Changes since the last full backup. & The last full backup plus only the most recent differential.\\
\bottomrule
\end{tabular}
\end{center}

There is a consistent trade-off. Smaller and faster backups produce longer or
more complex restores. A full backup is heavy to create but trivial to
restore; an incremental is light to create but may need many files chained
together at restore time; a differential sits in between.

\begin{example}
On Monday the office desktop is given a full backup. On Tuesday, a single
document changes. With an incremental strategy, Tuesday's backup contains
only that one change. With a differential strategy, Tuesday's backup contains
that change measured against Monday's full backup. If only that one file
changed, the two results look identical on Tuesday --- the difference appears
on Wednesday, when the incremental again holds only Wednesday's change, while
the differential holds Tuesday's change plus Wednesday's change.
\end{example}

\section{Restore scenarios}
Restore behaviour is the part candidates most often miss, because it is not
visible on the day of the backup --- only when recovery is needed.

\begin{example}[Restore from an incremental chain]
A full backup is taken Sunday. Incrementals are taken Monday, Tuesday and
Wednesday. On Thursday the disk fails. To recover Wednesday's data you need
the Sunday full backup and then \emph{all three} incrementals --- Monday,
Tuesday and Wednesday --- applied in order. Missing any one link loses the
changes it carried.
\end{example}

\begin{example}[Restore from a differential chain]
A full backup is taken Sunday. Differentials are taken Monday, Tuesday and
Wednesday. On Thursday the disk fails. To recover Wednesday's data you need
only the Sunday full backup and the \emph{most recent} differential ---
Wednesday's. The earlier differentials are not required, because Wednesday's
already includes all changes since Sunday.
\end{example}

Both scenarios produce the same recovered data. The difference is operational:
the incremental restore is longer and depends on every tape or file in the
chain; the differential restore is shorter and depends on only two sets.

\section{Image and file-level backups}
Beyond the three time-based types, two coverage-based terms appear in exams.

\begin{definitionbox}{Image backup (system image)}
A complete clone of a system --- operating system, settings, applications and
data --- that can restore a whole machine.
\end{definitionbox}

\begin{definitionbox}{File-level backup}
A backup of selected files and folders rather than a whole system.
\end{definitionbox}

An image restores a working machine in one step, which makes it ideal after a
drive failure or a corrupted operating system. A file-level backup is smaller
and more flexible for everyday documents, but recovering a machine from it
means reinstalling the operating system first and then restoring the files.
The office desktop could use an image for disaster recovery and file-level
backups for its working documents.

\section{The 3-2-1 rule, off-site and offline copies}
The \textbf{3-2-1 rule} is the standard exam reminder for resilient backup:
\begin{enumerate}
  \item \textbf{3} copies of the data --- the original plus two backups.
  \item On \textbf{2} different types of media (for example, an external disk
    and cloud storage).
  \item With \textbf{1} copy kept \textbf{off-site}, at a different location.
\end{enumerate}
The numbers map to copies, media and location, in that order. A common trap
is to swap them, or to read the ``2'' as two off-site copies.

An \textbf{off-site} copy is kept at a different physical place, so a single
fire, flood or theft cannot destroy every copy at once. An \textbf{offline}
copy is disconnected from the network; it is especially useful against
ransomware, which spreads through connected drives. A \textbf{cloud backup}
is an off-site copy held on remote servers. These are not mutually exclusive:
the office desktop's cloud backup is both off-site and, when not mounted, can
serve as an offline copy as well.

\begin{example}
The office desktop keeps its working files on the local SSD (copy 1), an
external USB disk (copy 2, different medium, kept offline), and a cloud
account (copy 3, off-site). This satisfies the 3-2-1 rule: three copies, two
media types, one off-site.
\end{example}

\section{Why RAID is not a backup}
RAID (\textbf{Redundant Array of Independent Disks}) combines several disks
for performance or redundancy. If one disk fails, a mirrored or parity RAID
keeps the system running. That protects \emph{uptime}, not \emph{recovery}.
RAID is not a backup because it does not protect against the things a backup
exists for: accidental deletion, corruption or ransomware. When a file is
deleted, or an encrypted copy is written by ransomware, the change is written
to every member of the RAID array at once. There is no earlier clean copy to
go back to. Redundancy duplicates the current state --- including the damage.

\section{Backup rotation}
\textbf{Backup rotation} means keeping several dated backup sets and reusing
them on a fixed schedule, commonly the grandfather--father--son scheme (daily,
weekly, monthly). Rotation matters because a problem such as silent
corruption may not be noticed for days; if the single backup set has already
been overwritten with the corrupted data, there is no clean copy left.
Keeping multiple versions across time ensures that an older, clean copy
remains available.

\subsection*{Exercises}
\begin{enumerate}[label=\arabic*.]
  \item Define full, incremental and differential backup in one line each.
  \item A full backup runs Sunday; differentials run Monday to Wednesday. The
    disk fails on Thursday. Which sets are needed to restore Wednesday's
    data?
  \item State the 3-2-1 rule and explain what each number refers to.
  \item Why is RAID not a substitute for a backup? Give one concrete danger.
  \item Explain the purpose of backup rotation.
\end{enumerate}

\subsection*{Solutions}
\begingroup\small
\begin{enumerate}[label=\arabic*.]
  \item Full: copies everything selected each run. Incremental: copies changes
    since the last backup of any type. Differential: copies changes since the
    last full backup.
  \item The Sunday full backup and the most recent (Wednesday) differential
    only. The earlier differentials are unnecessary because Wednesday's
    includes all changes since the full backup.
  \item Three copies of the data, on two different types of media, with one
    copy kept off-site.
  \item RAID provides redundancy for uptime, not recovery. A deletion,
    corruption or ransomware write is applied to every RAID member at once,
    so no earlier clean copy survives.
  \item Rotation retains several dated backup versions so that an older, clean
    copy is still available if corruption is discovered only after the newest
    backup has already been overwritten.
\end{enumerate}
\endgroup
\end{document}
